\begin{figure}[htbp]
    \centering
    \begin{minipage}{0.45\textwidth}
        \centering
        \begin{tikzpicture}[>=Stealth, thick, scale=0.9]

            % 1. Define Coordinates
            \coordinate (O) at (0,0);
            \coordinate (A_dir) at (100:6); % Direction of Va
            \coordinate (B_dir) at (10:6);  % Direction of Vb

            % Define the main radius
            \def\R{5}
            \coordinate (A) at (100:\R);
            \coordinate (B) at (10:\R);

            % Define P at an interpolated angle
            % Let's assume t is roughly 0.4 for visual similarity
            \coordinate (P) at (55:\R);

            % 2. Calculate Intersections for Parallelogram
            % We need A' on OA such that PA' // OB
            % We need B' on OB such that PB' // OA

            % Path for OA and OB
            \path[name path=rayOA] (O) -- (A_dir);
            \path[name path=rayOB] (O) -- (B_dir);

            % Path from P parallel to OB (to find A')
            \path[name path=lineFromP_parOB] (P) -- ++($(O)-(B_dir)$) -- ++($(B_dir)-(O)$);

            % Path from P parallel to OA (to find B')
            \path[name path=lineFromP_parOA] (P) -- ++($(O)-(A_dir)$) -- ++($(A_dir)-(O)$);

            % Find intersections
            \path [name intersections={of=rayOA and lineFromP_parOB, by=A_prime}];
            \path [name intersections={of=rayOB and lineFromP_parOA, by=B_prime}];

            % 3. Draw Geometry

            % Draw the main arc
            \draw (A) arc (100:10:\R);

            % Draw Vectors OA, OB, OP
            \draw[->] (O) -- (A) node[below left] {$\mathbf{v}_a$};
            \draw[->] (O) -- (B) node[below left] {$\mathbf{v}_b$};
            \draw[->] (O) -- (P) node[below left, xshift=-6pt, yshift=-2pt] {$\mathbf{v}_p$};

            % Draw Components (A' and B')
            % Vectors Va' and Vb' lie on the axes
            \draw[->, ultra thick, blue!80!black] (O) -- (A_prime) node[midway, left] {$\mathbf{v}_a'$};
            \draw[->, ultra thick, blue!80!black] (O) -- (B_prime) node[midway, below] {$\mathbf{v}_b'$};

            % Draw Dashed Parallel Lines
            \draw[dashed] (P) -- (A_prime);
            \draw[dashed] (P) -- (B_prime);

            % 4. Labels and Angles

            % Points labels
            \node[left] at (A_prime) {$A'$};
            \node[above right] at (B_prime) {$B'$};
            \node[above] at (A) {$A$};
            \node[right] at (B) {$B$};
            \node[above right] at (P) {$P$};

            % Angle arcs

            % Angle (1-t)theta between OA and OP
            \pic [draw, angle radius=1.0cm, pic text={$(1-t)\theta$}, angle eccentricity=1.6, font=\footnotesize] {angle = P--O--A};

            % Angle t*theta between OP and OB
            \pic [draw, angle radius=0.8cm, pic text={$t\theta$}, angle eccentricity=1.6, font=\footnotesize] {angle = B--O--P};

            % Internal angles (pi - theta) at A' and B'
            % Using small arcs inside the parallelogram
            \pic [draw, angle radius=0.6cm, pic text={$\pi-\theta$}, angle eccentricity=1.8, font=\footnotesize] {angle = O--A_prime--P};
            \pic [draw, angle radius=0.6cm, pic text={$\pi-\theta$}, angle eccentricity=1.8, font=\footnotesize] {angle = P--B_prime--O};

            % Angle t*theta at P (alternate interior) - optional based on diagram complexity
            % \pic [draw, angle radius=1.0cm, "$t\theta$"] {angle = A_prime--P--O};

        \end{tikzpicture}
    \end{minipage}
    \hfill
    \begin{minipage}{0.45\textwidth}
        \centering
        \begin{tikzpicture}[scale=2,every node/.style={font=\small},line width=0.8pt]
    \coordinate (Root) at (0,0);
    \coordinate (Internal) at (0, -0.3);
    \coordinate (A) at (-0.5, -2);
    \coordinate (B) at (0.5, -2);

    \draw (Root) -- (Internal);
    \draw (Internal) -| (A) node[pos=0.7, left] {$l_a$};
    \draw (Internal) -| (B) node[pos=0.7, right] {$l_b$};

    \node[below] at (A) {A};
    \node[below] at (B) {B};
    \node[above right] at (Internal) {P};
        \end{tikzpicture}
    \end{minipage}
    \caption{Geometric diagram of the ancestor state reconstruction algorithm. }
    \label{fig:algorithm-asr}
\end{figure}